\section{Related Work}
\label{sec:related}

\subsection{Centralized prioritized planning and conflict search}

Guo et al.'s ICARCV 2024 method is the baseline for this study. It formulates
large-scale multiple-AGV scheduling as prioritized assignment and route
planning, resolving conflicts against previously planned higher-priority
vehicles \cite{guo2024}. The sequential approach avoids joint search over all
vehicle paths, but its result depends on assignment and priority order.
Prioritized wait adjustment instead starts from individual paths and shifts
waiting using safe intervals; its preliminary evaluation reports 50--250
agents and trades solution quality for runtime \cite{andreychuk2018}.

Conflict-based search (CBS) offers a coupled alternative. Continuous-time CBS
methods model nonuniform action durations and geometric collisions
\cite{andreychuk2022}. He et al. extend CBS to multi-floor MAPF with
single-agent elevator occupancy and elevator reset intervals, explicitly
including lift conflicts in search \cite{he2026}. These methods provide
stronger conflict reasoning, while our experiments focus on job assignment,
prioritized route windows, and lift arbitration in a decomposed planner.

\subsection{Distributed multi-robot and AGV coordination}

Distributed manufacturing schedulers assign decisions to resource agents.
Zhou et al. use collaborative reinforcement-learning schedulers for
heterogeneous smart-factory units and consider rush orders and machine
failures \cite{zhou2021}. Maoudj et al. formulate a capacitated multi-AGV
scheduling problem with conflicting products and compare a decentralized
multi-agent approach with a mixed-integer model \cite{maoudj2023}. Draganjac
et al. present decentralized traffic management based on local coordination
for free-ranging industrial vehicles \cite{draganjac2020}. These works
address decentralized decisions or online manufacturing dynamics; they do
not evaluate the same combination of a central lift calendar and floor-local
route-leg reservations used here.

\subsection{Hierarchical and zone-based planning}

Spatial decomposition reduces the planning scope by assigning agents to
regions or layers. The spare-zone method of Zhao et al. uses a hierarchical
allocation layer and local coordination to resolve collisions and deadlocks
for vehicles following fixed paths \cite{zhao2021}. Pratissoli et al. use a
centralized, decoupled, multilayer traffic manager with topological sectors,
traffic-aware route planning, and time-expanded coordination in an industrial
setting \cite{pratissoli2024}. HMAPP combines high-level region sequences
with regional conflict-based searches \cite{hmapp2021}; hierarchical
large-scale multirobot replanning similarly combines global congestion-aware
routing and parallel low-level planning \cite{pan2024}. Verma et al.'s
dynamic resource reservation method allows multiple reservations at shared
resource points and can be implemented centrally or with local controllers
\cite{verma2024}. These approaches motivate layered coordination and local
waiting or conflict-resolution areas. The present implementation instead
partitions local route windows by floor while retaining central job
assignment and lift booking.

\subsection{Learning-based and optimization-based AGV scheduling}

Scheduling quality and traffic efficiency have also been addressed with
optimization, learning, and large-neighborhood methods. Wang et al. use a
neural network to predict travel time in a multi-state AGV scheduling method
for flexible manufacturing \cite{wang2022}. Lu et al. combine workstation
assignment rules with a space-time scheduling model and cooperative
optimization for conflict-free AGV routes \cite{lu2023}. Hu et al. combine
reservation timetables, task allocation, and variable-neighborhood search
for multi-load AGVs; their experiments include fleets from 3 to 200 vehicles
\cite{hu2022}. MAPF-LNS2 repairs collisions by repeatedly replanning subsets
of agents and reports results up to 8,000 agents on a congested warehouse
benchmark under a one-hour limit \cite{mapflns2022}. These approaches target
search quality, assignment, or scale, but the cited formulations do not
evaluate our specific central-lift and floor-server split.

\subsection{Multi-floor warehouses and lifts}

Elevator-aware warehouse routing is studied directly by Ma et al. for a
double-warehouse system connected by elevators \cite{ma2014}. Werner analyzes
MAPF in high-rack warehouses where agents remain on their assigned levels
and elevators transport goods vertically; the method constructs waiting
zones for elevator use \cite{werner2025}. He et al.'s CBS-E formulation
models elevator capacity and reset-time conflicts \cite{he2026}. These
papers demonstrate that lift state and waiting near elevator shafts matter
to multi-floor routing. Our evaluated contribution combines this resource
with prioritized job assignment, floor-local route-leg planning, and
central lift booking. The claim is limited to this implemented combination
and the experiments in this paper; we make no claim of priority over all
related systems.

\subsection{Contributions}

This work reimplements the ICARCV 2024 prioritized-planning baseline with
corrections to the original demonstration, then evaluates a hierarchical
central-plus-floor-server extension. The central component assigns jobs and
books lifts, floor planners resolve local route conflicts, and lift-lobby
waiting points support delays near floor transitions. An independent checker
validates completed schedules. Results compare centralized, FCFS, and hybrid
variants over seeded job sets, fleet sizes, job scales, and parameter settings
for the tested factory layouts \cite{projectRepo}.
